\begin{dang}{Âm và nhạc âm}
\begin{itemize}
\item Tần số âm cơ bản do dây đàn phát ra: $f_0=\dfrac{v}{2\ell}$
\item Hoạ âm bậc k của dây đàn có tần số: $f_n=kf_0$
\end{itemize}
\begin{itemize}
\item Nếu dùng âm thoa để kích thích dao động một cột khí (chiều
cao cột khí có thể thay đổi bằng cách thay đổi mực nước), khi có sóng dừng trong cột khí thì đầu kín luôn luôn là nút, còn đầu hở có thể nút
hoặc bụng.
\begin{itemize}
\item Nếu đầu hở là bụng thì âm nghe được là to nhất. Chiều dài cột khí:  
\begin{align*}
\ell=\left(2k-1\right).\dfrac{\lambda}{4}
\Rightarrow \ell_{\min}=\dfrac{\lambda}{4}
\end{align*}
\item Nếu đầu hở là nút thì âm nghe được là nhỏ nhất. Chiều dài cột khí: 
\begin{align*}
\ell=k.\dfrac{\lambda}{2}\Rightarrow \ell_{\min}=\dfrac{\lambda}{2}
\end{align*}
\end{itemize}
\Note{\begin{itemize}
\item Khi âm phát ra to nhất thì đầu hở là bụng sóng.
\item Khi gõ vào thành ống tần số phát ra là tần số cơ bản.
\end{itemize}
}
\end{itemize}
\end{dang}